\section*{Chapitre \ref{ch:g12:matrix} --- Matrices et graphes}

\begin{solution}{exo:g12:matrix:1}
\[
A + B = \begin{pmatrix} 3 & 2\\ 1 & 2\end{pmatrix}, \quad
AB = \begin{pmatrix} 4 & 2\\ 1 & 1\end{pmatrix}, \quad
BA = \begin{pmatrix} 2 & 4\\ 1 & 3\end{pmatrix}, \quad
A^2 = \begin{pmatrix} 1 & 4\\ 0 & 1\end{pmatrix}.
\]
Noter $AB \neq BA$.
\end{solution}

\begin{solution}{exo:g12:matrix:2}
$\det A = 6 - 5 = 1 \neq 0$~:
$A^{-1} = \begin{pmatrix} 3 & -5\\ -1 & 2 \end{pmatrix}$.
$\det B = 12 - 12 = 0$~: $B$ n'est pas \omterm{def:g12:matrix:inverse}{inversible}.
\end{solution}

\begin{solution}{exo:g12:matrix:3}
Le système est $AX = Y$ avec $A$ comme dans l'\cref{exo:g12:matrix:2} et
$Y = \begin{pmatrix} 1\\ 2\end{pmatrix}$~:
\[
X = A^{-1}Y = \begin{pmatrix} 3 & -5\\ -1 & 2\end{pmatrix}
\begin{pmatrix} 1\\ 2\end{pmatrix}
= \begin{pmatrix} -7\\ 3\end{pmatrix}:
\qquad x = -7,\ y = 3 .
\]
\end{solution}

\begin{solution}{exo:g12:matrix:4}
\emph{1.} $N^2 = \begin{pmatrix} 0&1\\0&0\end{pmatrix}
\begin{pmatrix} 0&1\\0&0\end{pmatrix} = 0$, et $I$ commute avec toute
\omterm{def:g12:matrix:matrix}{matrice}.

\emph{2.} Comme les deux termes commutent, la formule du binôme
s'applique et tous les termes contenant $N^2$ s'annulent~:
\[
A^k = (2I + N)^k = 2^k I + k\,2^{k-1} N
= \begin{pmatrix} 2^k & k\,2^{k-1}\\ 0 & 2^k \end{pmatrix}.
\]
Vérification pour $k = 2$~:
$A^2 = \begin{pmatrix} 2&1\\0&2\end{pmatrix}^2
= \begin{pmatrix} 4&4\\0&4\end{pmatrix}$, et la formule donne
$2^2 = 4$, $2 \times 2 = 4$. \checkmark
\end{solution}

\begin{solution}{exo:g12:matrix:5}
\emph{Récurrence.} Pour $n = 1$~:
$A^1 = \begin{pmatrix} 0&1\\1&1\end{pmatrix}
= \begin{pmatrix} F_0 & F_1\\ F_1 & F_2\end{pmatrix}$. Supposer la
formule pour $n$~; alors
\[
A^{n+1} = A^n A
= \begin{pmatrix} F_{n-1} & F_n\\ F_n & F_{n+1}\end{pmatrix}
\begin{pmatrix} 0 & 1\\ 1 & 1\end{pmatrix}
= \begin{pmatrix} F_n & F_{n-1} + F_n\\ F_{n+1} & F_n + F_{n+1}\end{pmatrix}
= \begin{pmatrix} F_n & F_{n+1}\\ F_{n+1} & F_{n+2}\end{pmatrix}.
\]
\emph{Identité.} Pour les \omterm{def:g12:matrix:matrix}{matrices} $2\times2$, le développement montre
$\det(MN) = \det M \det N$~; d'où
$\det(A^n) = (\det A)^n = (-1)^n$, et
$\det A^n = F_{n-1}F_{n+1} - F_n^2$. (C'est l'\emph{identité de Cassini}.)
\end{solution}

\begin{solution}{exo:g12:matrix:6}
\emph{1.} En ordonnant les sommets $1, 2, 3$~:
\[
M = \begin{pmatrix} 0&1&1\\ 0&0&1\\ 1&0&0\end{pmatrix}, \quad
M^2 = \begin{pmatrix} 1&0&1\\ 1&0&0\\ 0&1&1\end{pmatrix}, \quad
M^3 = \begin{pmatrix} 1&1&1\\ 1&0&1\\ 1&0&1 \end{pmatrix}.
\]

\emph{2.} $\bigl(M^3\bigr)_{11} = 1$~: exactement un chemin fermé de
longueur $3$ au sommet $1$, à savoir $1 \to 2 \to 3 \to 1$. (Le chemin
$1 \to 3 \to 1$ n'a que longueur $2$, et $1 \to 3$ puis $3\to1$ puis
$1\to3$ se termine en $3$.)
\end{solution}

\begin{solution}{exo:g12:matrix:7}
\emph{1.} $a_{n+1} = 0.8a_n + 0.3b_n$, $b_{n+1} = 0.2a_n + 0.7b_n$~:
$M = \begin{pmatrix} 0.8 & 0.3\\ 0.2 & 0.7\end{pmatrix}$.

\emph{2.} $MX = X$ donne $0.8a + 0.3b = a$, \emph{c.-à-d.} $0.3b = 0.2a$,
donc $b = \frac23 a$~; avec $a + b = 1$~: $a = 0.6$, $b = 0.4$.

\emph{3.} En utilisant $b_n = 1 - a_n$~:
$a_{n+1} = 0.8a_n + 0.3(1 - a_n) = 0.5a_n + 0.3$, donc
\[
c_{n+1} = a_{n+1} - 0.6 = 0.5a_n + 0.3 - 0.6 = 0.5(a_n - 0.6) = 0.5\,c_n .
\]
D'où $c_n = 0.5^n c_0 \to 0$~: $a_n \to 0.6$ et $b_n \to 0.4$, quelle que
soit la répartition initiale.
\end{solution}

\begin{solution}{exo:g12:matrix:8}
\emph{1.} $\det P = -2$, donc
$P^{-1} = -\frac12\begin{pmatrix} -1 & -1\\ -1 & 1\end{pmatrix}
= \frac12\begin{pmatrix} 1 & 1\\ 1 & -1\end{pmatrix}$. Puis
\[
AP = \begin{pmatrix} 4 & 2\\ 4 & -2 \end{pmatrix}, \qquad
D = P^{-1}AP = \frac12\begin{pmatrix} 1&1\\1&-1\end{pmatrix}
\begin{pmatrix} 4&2\\4&-2\end{pmatrix}
= \begin{pmatrix} 4 & 0\\ 0 & 2\end{pmatrix}.
\]

\emph{2.} De $A = PDP^{-1}$, une récurrence immédiate donne
$A^n = PD^nP^{-1}$ avec
$D^n = \begin{pmatrix} 4^n & 0\\ 0 & 2^n\end{pmatrix}$, donc
\[
A^n = P D^n P^{-1}
= \begin{pmatrix} 4^n & 2^n\\ 4^n & -2^n\end{pmatrix}\cdot
\frac12\begin{pmatrix} 1&1\\1&-1\end{pmatrix}
= \frac12\begin{pmatrix} 4^n + 2^n & 4^n - 2^n\\
4^n - 2^n & 4^n + 2^n\end{pmatrix}.
\]
Appliquer $A^n$ à $X_0 = \begin{pmatrix} u_0\\v_0\end{pmatrix}$
reproduit exactement les formules de l'\cref{ex:g12:matrix:coupled}.
\end{solution}

\begin{solution}{pb:g12:matrix:1}
\textbf{1.} $AB = \begin{pmatrix} 2 & 1\\ 4 & 3\end{pmatrix}$ et
$BA = \begin{pmatrix} 3 & 4\\ 1 & 2\end{pmatrix}$~: la
multiplication des \omterm{def:g12:matrix:matrix}{matrices} n'est pas commutative --- $B$ échange
les colonnes à droite, les lignes à gauche.

\textbf{2.} \omterm{def:g11:vect:det}{Déterminant} $6 - 5 = 1$, d'où l'inverse
$\begin{pmatrix} 3 & -1\\ -5 & 2\end{pmatrix}$. En l'appliquant à
$\binom{4}{7}$~: $x = 12 - 7 = 5$, $y = -20 + 14 = -6$.

\textbf{3.} $N^2 = 0$. Alors $(I + N)^n = I + nN$ par récurrence~:
$(I + nN)(I + N) = I + (n+1)N + nN^2 = I + (n+1)N$.

\textbf{4.} $A = \begin{pmatrix} 0&1&1\\ 1&0&1\\ 1&1&0
\end{pmatrix}$, et $A^3$ a pour coefficients diagonaux $2$~: depuis
chaque sommet, exactement deux chemins fermés de longueur $3$ (le
triangle parcouru dans un sens ou dans l'autre) --- le théorème de
dénombrement à l'œuvre.

\textbf{5.} $D^n = \begin{pmatrix} 2^n & 0\\ 0 & 2^{-n}
\end{pmatrix}$~: une direction explose, l'autre s'éteint --- sur la
diagonale, les destins sont des \omterm{def:g11:seq:geometric}{suites géométriques} indépendantes.

\textbf{6.} $F^2 = \begin{pmatrix} 2 & 1\\ 1 & 1\end{pmatrix}$,
$F^3 = \begin{pmatrix} 3 & 2\\ 2 & 1\end{pmatrix}$,
$F^4 = \begin{pmatrix} 5 & 3\\ 3 & 2\end{pmatrix}$~: du Fibonacci
partout, d'où la conjecture énoncée.

\textbf{7.} Si $F^n = \begin{pmatrix} F_{n+1} & F_n\\ F_n &
F_{n-1}\end{pmatrix}$, alors
\[
F^{n+1} = F^n F =
\begin{pmatrix} F_{n+1} + F_n & F_{n+1}\\
F_n + F_{n-1} & F_n \end{pmatrix}
= \begin{pmatrix} F_{n+2} & F_{n+1}\\ F_{n+1} & F_n
\end{pmatrix} :
\]
hérédité~; l'initialisation en $n = 1$ n'est autre que $F$
lui-même, avec la convention $F_0 = 0$ (qui prolonge la récurrence
vers l'arrière).

\textbf{8.} $\det F = -1$, donc
$\det(F^n) = (\det F)^n = (-1)^n$~; et directement
$\det(F^n) = F_{n+1}F_{n-1} - F_n^2$~: c'est Cassini, en une ligne.
(La règle du produit pour les \omterm{def:g11:vect:det}{déterminants} $2 \times 2$ est un
agréable développement de cinq minutes.)

\textbf{9.} Coefficient en haut à droite de $F^m F^n$~:
$F_{m+1}F_n + F_m F_{n-1}$~; celui de $F^{m+n}$~: $F_{m+n}$. Pour
$m = n = 3$~: $F_4 F_3 + F_3 F_2 = 3 \times 2 + 2 \times 1 = 8 =
F_6$.

\textbf{10.} Pour $k = 1$~: c'est immédiat. Si $F_n \mid F_{kn}$,
la formule d'addition avec $m = kn$ donne
$F_{(k+1)n} = F_{kn+1}F_n + F_{kn}F_{n-1}$~: les deux termes sont
des multiples de $F_n$. Donc $F_n \mid F_{kn}$ pour tout $k$~:
vérification, $F_3 = 2$ divise $F_6 = 8$ et $F_9 = 34$.

\textbf{11.} $F^{100} = F^{64} F^{32} F^4$~: sept élévations au
carré ($F^2, F^4, \dots, F^{64}$) plus deux combinaisons --- neuf
multiplications au lieu de quatre-vingt-dix-neuf. C'est l'astuce de
la table de doublements des scribes égyptiens, transposée des
nombres aux \omterm{def:g12:matrix:matrix}{matrices}~: écrire $100$ en binaire, puis multiplier les
doublements utiles.

\textbf{12.} $0.8 + 0.2 = 1$ et $0.4 + 0.6 = 1$~: demain il fera
\emph{un} temps ou un autre --- chaque colonne est une \omterm{def:g11:prob:rv}{loi} de
\omterm{def:g10:proba:distribution}{probabilité} complète, si bien que les \omterm{def:g10:proba:distribution}{probabilités} se conservent.

\textbf{13.} Demain~: $\binom{0.8}{0.2}$. Après-demain~:
$M\binom{0.8}{0.2} = \binom{0.72}{0.28}$.

\textbf{14.} $Mv = v$ avec $v = \binom{s}{r}$ et $s + r = 1$~:
$0.8s + 0.4r = s$ donne $0.4r = 0.2s$, soit $s = 2r$, d'où
$v = \binom{2/3}{1/3}$. À long terme, deux jours sur trois sont
ensoleillés --- quel que soit le temps qu'il fait aujourd'hui.

\textbf{15.} À partir de $\binom01$~: composantes ensoleillées
$0.4$, $0.56$, $0.624$, $0.6496$~; écarts à $\frac23$~: $0.267$,
$0.107$, $0.043$, $0.017$ --- chaque pas multiplie l'écart par
exactement $0.4$ (la seconde valeur propre de la machine)~:
convergence géométrique vers l'état stationnaire.

\textbf{16.} Colonnes (issues de $A$, $B$, $C$)~:
$P = \begin{pmatrix} 0 & 0 & 1\\ \frac12 & 0 & 0\\
\frac12 & 1 & 0\end{pmatrix}$. État stationnaire~: $v_A = v_C$,
$v_B = \frac{v_A}{2}$, $v_C = \frac{v_A}{2} + v_B$~; en sommant à
$1$~: $v = \left(\frac25, \frac15, \frac25\right)$. Classement~:
$A$ et $C$ à égalité en tête, $B$ dernière.

\textbf{17.} $C$ reçoit \emph{tout} le trafic de $B$ et la moitié
de celui de $A$, et il renvoie tout vers $A$~: l'état stationnaire
mesure où l'internaute \emph{passe du temps}, et non combien de
liens pointent vers la page --- un lien depuis une page fréquentée
pèse plus que plusieurs liens depuis des pages désertes. Cette
pondération récursive est précisément l'idée fondatrice du moteur~;
l'amortissement gère les pièges et les impasses.

\textbf{18.} $A^n = P D^n P^{-1}$ donne
$u_n = \frac{4^n + 2^n}{2}$ (et $v_n = \frac{4^n - 2^n}{2}$).
Vérification~: $u_1 = 3$, $u_2 = 10$, $u_3 = 36$~; directement,
$(1,0) \to (3,1) \to (10,6) \to (36, 28)$~: cela concorde.

\textbf{19.} La diagonalisation change de \omterm{def:g10:coordgeom:system}{coordonnées} pour des
\omterm{def:g10:coordgeom:system}{coordonnées} dans lesquelles le système couplé se disloque en \omterm{def:g11:seq:geometric}{suites
géométriques} indépendantes --- chaque valeur propre court sa propre
course. L'état stationnaire de Markov est le \omterm{def:g11:vect:vector}{vecteur} propre associé
à la valeur propre $1$, et la vitesse de convergence de la question
15 est la valeur propre suivante~: la machine à météo était depuis
le début une histoire d'éléments propres.

\textbf{20.} Tenue de comptes~: un système est une seule \omterm{def:g10:algebra:equation}{équation}
matricielle, résolue par un seul inverse. Dénombrement~: les
puissances de la \omterm{def:g12:matrix:graph}{matrice d'adjacence} comptent les chemins, les liens
et les connexions. Évolution~: les puissances de la machine
conduisent les états vers leurs destins, et les directions propres
(la direction dorée de Fibonacci, l'état stationnaire de la météo,
le \omterm{def:g11:vect:vector}{vecteur} de classement du web) \emph{sont} ces destins. L'algèbre
linéaire, dans les volumes universitaires, est la science de tout
cela.
\end{solution}
